\begin{theorem}[Trace-lattice formulas]\label{mod:ramification-traceideals}
Let $K/F$ be a finite Galois valuative extension of nonarchimedean local
fields, put
\[
 e=e(K/F),\qquad D=d_{K/F}=v_K(\mathfrak D_{K/F}),
\]
and identify $D$ with the nonnegative integer
\texttt{differentExponent} constructed in
\ref{mod:ramification-different}.  Lean keeps the following two hypotheses
on a chosen $\pi\in\mathcal O_K$ explicit:
\[
 \operatorname{ord}_K(\pi)=1,
 \qquad
 \operatorname{Algebra.adjoin}_{\mathcal O_F}\{\pi\}=\mathcal O_K.
\]
The second equality first implies $F(\pi)=K$.  The derivative of the
minimal polynomial of $\pi$ is then evaluated as
\[
 f_\pi'(\pi)
 =\prod_{\substack{\sigma\in\operatorname{Gal}(K/F)\\\sigma\ne1}}
   (\pi-\sigma\pi),
\]
so its upper normalized order is exactly $D$ by the Hilbert-sum theorem of
\ref{mod:ramification-different}.  Mathlib's monogenic trace-duality theorem
therefore identifies the genuine integral trace dual with
$\mathfrak p_K^{-D}$; the desired lattice identity is not assumed through a
typeclass.

For every $a\in\mathbf Z$, including all negative depths, Lean proves
\[
 \boxed{
 \left(\mathfrak p_K^a\right)^\vee
 =\{z\in K:\operatorname{Tr}_{K/F}(z\mathfrak p_K^a)
                    \subseteq\mathcal O_F\}
 =\mathfrak p_K^{-a-D}.}
\]
Multiplication of integer-indexed lattices transports the integral result to
arbitrary $a$ and fixes both the minus sign on $a$ and the negative
different shift.

Lean then proves the exact trace image, not merely its containment:
\[
 \boxed{
 \operatorname{Tr}_{K/F}(\mathfrak p_K^a)
 =\mathfrak p_F^{\left\lfloor(a+D)/e\right\rfloor}}
 \qquad(a\in\mathbf Z).
\]
Here division in the Lean exponent is Euclidean integer division by the
positive integer $e$, hence is the displayed floor even when $a+D<0$.
The theorem is stated both as equality of the mapped submodule and in the
elementwise form
\[
 y\in\mathfrak p_F^{\left\lfloor(a+D)/e\right\rfloor}
 \quad\Longleftrightarrow\quad
 \exists x\in\mathfrak p_K^a,\ \operatorname{Tr}_{K/F}(x)=y.
\]
Thus the zero target and its zero preimage are included.  The reverse
inclusion is proved from trace duality: if every trace lay one layer deeper,
an element of base order $-\lfloor(a+D)/e\rfloor-1$ would contradict the
dual-lattice formula.  Scaling one resulting exact-depth trace then supplies
every target element; trace surjectivity is never an input hypothesis.

The annihilator exponent used in the manuscript is recorded with its exact
endpoint convention:
\[
 \left\lceil\frac{-a-D}{e}\right\rceil
 =\left\lfloor\frac{-(a+D)+e-1}{e}\right\rfloor
 =-\left\lfloor\frac{a+D}{e}\right\rfloor.
\]
Consequently, for every additive character $\psi$ of $F$ and every integer
$a$, triviality of $\psi\circ\operatorname{Tr}_{K/F}$ on
$\mathfrak p_K^a$ is equivalent to triviality of $\psi$ on the displayed
exact trace-image lattice.  With the convention that conductor $n$ means
that $\mathfrak p_F^{-n}$ is the largest triviality lattice, Lean proves
\[
 \operatorname{IsAdditiveConductor}_F(\psi,n)
 \Longrightarrow
 \operatorname{IsAdditiveConductor}_K
   (\psi\circ\operatorname{Tr}_{K/F},e n+D).
\]
In the manuscript's numerical conductor notation, this says
\[
 n_K(\psi\circ\operatorname{Tr}_{K/F})=e\,n_F(\psi)+D.
\]
For the totally ramified prime-degree situation of
\ref{cor:additive-conductor}, this is exactly
$\ell n_F(\psi)+d_{K/F}$.

Finally assume
\[
 e=2,\qquad [K:F]=2,\qquad r\geq0,
 \qquad r+D=2q+1.
\]
Lean defines trace on quotient representatives and proves the representative
formula
\[
 [x]\longmapsto[\operatorname{Tr}_{K/F}(x)]
 \colon
 \mathfrak p_K^r/\mathfrak p_K^{r+1}
 \longrightarrow
 \mathfrak p_F^q/\mathfrak p_F^{q+1}.
\]
It derives residue degree one from $e=[K:F]=2$, rather than postulating
surjectivity of the residue map, and proves the exact kernel statement
\[
 x\in\mathfrak p_K^r
 \Longrightarrow
 \bigl(\operatorname{Tr}_{K/F}(x)\in\mathfrak p_F^{q+1}
       \iff x\in\mathfrak p_K^{r+1}\bigr).
\]
Together with the exact trace-image formula this makes the displayed map an
$\mathcal O_F$-linear equivalence.  Equivalently, for
$x\in\mathfrak p_K^r$,
\[
 \operatorname{ord}_F(\operatorname{Tr}_{K/F}x)=q
 \iff \operatorname{ord}_K(x)=r;
\]
in particular, zero belongs to neither exact shell.  In the prime-cyclic
wild quadratic specialization, Lean substitutes the proved equality
$D=t+1$ and exposes the same equivalence with that shift.

The two endpoints of Lemma~\ref{lem:quadratic-graded-trace} are also formal
theorems.  If $D=2d_\tau+1$, then
\[
 \operatorname{ord}_F(2)=d_\tau.
\]
For every integral unit $w$ and every $\widetilde{\bar w}\in\mathcal O_F$
whose residue maps to that of $w$,
\[
 \operatorname{Tr}_{K/F}(w)-2\widetilde{\bar w}
 \in\mathfrak p_F^{d_\tau+1}.
\]
Writing $T=D$, if $m=2d+1>T$ and $r=m-T$, with the strict inequality retained, Lean
constructs the equivalence at depths $r$ and $d$ and proves the downstream
nonzero-shell form
\[
 \operatorname{ord}_F(-\operatorname{Tr}_{K/F}x)=d
 \iff \operatorname{ord}_K(x)=r
 \qquad(x\in\mathfrak p_K^r).
\]

\LeanFile{LanglandsFirstMainLemma/Ramification/TraceIdeals.lean}
\lean{LanglandsFirstMainLemma.mem_traceDual_iff_trace_mul_mem_lattice_zero}
\lean{LanglandsFirstMainLemma.traceDual_lattice_zero}
\lean{LanglandsFirstMainLemma.traceDual_lattice}
\lean{LanglandsFirstMainLemma.trace_lattice_image_eq}
\lean{LanglandsFirstMainLemma.mem_trace_lattice_image_iff}
\lean{LanglandsFirstMainLemma.trace_mem_lattice_floor}
\lean{LanglandsFirstMainLemma.trace_annihilator_ceiling_eq}
\lean{LanglandsFirstMainLemma.addCharTrivialOnLattice_compTrace_iff}
\lean{LanglandsFirstMainLemma.isAdditiveConductor_compTrace}
\lean{LanglandsFirstMainLemma.traceLatticeGraded}
\lean{LanglandsFirstMainLemma.traceLatticeGraded_mk}
\lean{LanglandsFirstMainLemma.quadraticGradedTrace}
\lean{LanglandsFirstMainLemma.quadraticGradedTrace_mk}
\lean{LanglandsFirstMainLemma.quadratic_trace_mem_succ_iff}
\lean{LanglandsFirstMainLemma.quadraticGradedTrace_bijective}
\lean{LanglandsFirstMainLemma.quadraticGradedTraceEquiv}
\lean{LanglandsFirstMainLemma.wildQuadraticGradedTraceEquiv}
\lean{LanglandsFirstMainLemma.quadratic_trace_shell_iff}
\lean{LanglandsFirstMainLemma.quadratic_ord_two}
\lean{LanglandsFirstMainLemma.quadratic_unit_trace_congr}
\lean{LanglandsFirstMainLemma.quadraticOddConductorGradedTraceEquiv}
\lean{LanglandsFirstMainLemma.quadratic_odd_conductor_neg_trace_shell_iff}
\leanok
\SourceText{Theorem~\ref{thm:norm-filtration}, part (b),
Corollary~\ref{cor:additive-conductor}, and
Lemma~\ref{lem:quadratic-graded-trace}.}
\uses{mod:ramification-different,mod:localfield-lattices,mod:localfield-extension,
mod:basic-characterconductors}
\FormalizationContract The different is connected to Mathlib's actual trace
dual through the chosen generating uniformizer and its minimal-polynomial
derivative.  Duality and exact trace images are proved at every integer
depth, with positive ramification index, Euclidean floor division, the
ceiling annihilator identity, and the additive-conductor sign all explicit.
No trace-surjectivity or trace-lattice assertion is assumed or stored in a
helper typeclass.  The quadratic API contains the representative map, exact
next-layer kernel, linear equivalence, exact-shell criterion, odd-different
unit endpoint, and strict $m>T$ specialization.
\end{theorem}
